\section{Requirements Specification}
\label{sec:requirements}

\subsection{Stakeholders and Personas}

Three personas drive requirement priority. They are listed in order of
decision-making weight for v1.0.

\begin{itemize}
\item \textbf{P1 --- The Working Photographer.} Shoots both film and digital.
Owns a scanner or pays for lab scans. Can name the difference between Portra
400 and Gold 200. Wants the iPhone frames in a project to intercut with the
film frames. Will not tolerate a result that ``looks like a filter'' and will
detect a wrong grain structure at 100\% magnification immediately.
\item \textbf{P2 --- The Filmmaker / Colorist.} Uses DaVinci Resolve
professionally. Understands printer lights, print emulation, and log
workflows. Wants the mobile result to be a legitimate starting point that can
be handed to a desktop timeline. Judges the product by whether its vocabulary
is used correctly.
\item \textbf{P3 --- The Serious Enthusiast.} Does not know what a shoulder is
but knows what looks good and is willing to learn. Is the volume of the
business. Requires that the default path from capture to good result be short,
and that the vocabulary be teachable rather than merely present.
\end{itemize}

The product tension across these personas is the central UX problem, and the
resolution is specified in Section~\ref{sec:ui}: a two-tier control surface in
which every P3 control is a named point on a P1/P2 parameter manifold, not a
separate simplified code path.

\subsection{Functional Requirements}

\begin{table}[htbp]
\caption{Functional Requirements, v1.0}
\label{tab:fr}
\begin{center}
\renewcommand{\arraystretch}{1.12}
\begin{tabularx}{\columnwidth}{@{}l>{\raggedright\arraybackslash}X c@{}}
\toprule
\textbf{ID} & \textbf{Requirement} & \textbf{Pri} \\
\midrule
FR-01 & Capture Apple ProRAW (DNG) at maximum supported resolution with user-selectable ISO, shutter, focus, and white balance. & MUST \\
FR-02 & Import DNG, Apple ProRAW, and 16-bit TIFF from Photos and Files. & MUST \\
FR-03 & Decode RAW to scene-referred linear working space with no vendor tone curve applied. & MUST \\
FR-04 & Provide $\geq 10$ negative stock profiles at launch spanning color negative, transparency, and monochrome. & MUST \\
FR-05 & Provide $\geq 4$ print stock profiles plus a bypass (scan-look) mode. & MUST \\
FR-06 & Expose development controls: push/pull in $1/3$ stop increments over $[-2,+3]$, contrast index, and agitation/temperature proxies. & MUST \\
FR-07 & Expose printer lights per channel in integer points over $[-12,+12]$ with a linked master. & MUST \\
FR-08 & Synthesize grain procedurally with per-stock structure, ISO dependence, and resolution-aware scaling. & MUST \\
FR-09 & Render halation with threshold, radius, intensity, and channel bias controls. & MUST \\
FR-10 & Provide optical simulation: diffusion, vignetting, chromatic aberration, bloom, and distortion. & SHOULD \\
FR-11 & Provide aging effects: dust, scratches, edge fog, expired-stock shift, light leaks. & SHOULD \\
FR-12 & All edits non-destructive; original RAW never modified. & MUST \\
FR-13 & Full edit history with undo/redo to arbitrary depth within a session and at least 50 steps persisted. & MUST \\
FR-14 & Save, name, and apply recipes (complete parameter graphs) across images. & MUST \\
FR-15 & Copy/paste individual stage parameters between images. & SHOULD \\
FR-16 & Export to HEIC, JPEG, 16-bit TIFF, and 16-bit PNG with selectable output color space. & MUST \\
FR-17 & Export a sidecar recipe file (JSON) and import same. & SHOULD \\
FR-18 & Batch apply a recipe to a selection of images. & SHOULD \\
FR-19 & Before/after comparison with split, toggle, and pin-to-stage modes. & MUST \\
FR-20 & Pipeline inspector showing per-stage input and output previews. & SHOULD \\
FR-21 & Preserve and write capture metadata (EXIF) plus a custom development-parameter block into exports. & SHOULD \\
FR-22 & Histogram in both scene-referred log and output display domains. & SHOULD \\
\bottomrule
\end{tabularx}
\end{center}
\end{table}

\subsection{Non-Functional Requirements}

\begin{table}[htbp]
\caption{Non-Functional Requirements}
\label{tab:nfr}
\begin{center}
\renewcommand{\arraystretch}{1.12}
\begin{tabularx}{\columnwidth}{@{}l>{\raggedright\arraybackslash}X@{}}
\toprule
\textbf{ID} & \textbf{Requirement} \\
\midrule
NFR-01 & Parameter change to updated preview: $\leq 16$\,ms at preview resolution on A17 Pro or later; $\leq 33$\,ms on A15. \\
NFR-02 & Full-resolution 48\,MP render to export: $\leq 2.5$\,s on A17 Pro. \\
NFR-03 & Peak GPU memory during full-resolution render: $\leq 900$\,MB. \\
NFR-04 & Application cold launch to camera-ready: $\leq 1.2$\,s. \\
NFR-05 & Capture shutter latency: $\leq 250$\,ms to shutter confirmation. \\
NFR-06 & No visible banding: quantization error $\leq 1$ code value at 16-bit output for any monotone ramp. \\
NFR-07 & Preview and export must be perceptually identical: $\Delta E_{00} \leq 1.0$ mean, $\leq 2.0$ 99th percentile, excluding grain. \\
NFR-08 & Grain must be deterministic given a seed: identical parameters yield identical output. \\
NFR-09 & Thermal: sustained editing for 10 minutes must not trigger severe thermal state on a device at 25\textdegree C ambient. \\
NFR-10 & Crash-free session rate $\geq 99.7\%$ at GA. \\
NFR-11 & Accessibility: all controls reachable by VoiceOver with meaningful labels; Dynamic Type support to XXL in all text UI. \\
NFR-12 & Storage: application-managed cache never exceeds 4\,GB and is evictable under system pressure. \\
NFR-13 & Offline-first: no network required for any v1.0 feature. \\
\bottomrule
\end{tabularx}
\end{center}
\end{table}

\subsection{Constraints}

\begin{itemize}
\item \textbf{CR-01.} iOS 17.0 minimum deployment target; devices with A14
Bionic or later. This is driven by Metal 3 feature availability, specifically
mesh-free tile shading and the fast resource heap aliasing the render graph
depends on.
\item \textbf{CR-02.} All processing on device. No image data leaves the
device in v1.0. This is both a privacy commitment and a latency requirement.
\item \textbf{CR-03.} No use of private API. The application must be
App Store distributable without entitlement exceptions.
\item \textbf{CR-04.} Film and print stock names are used descriptively.
Profiles are derived from published sensitometric data and from measurements of
physical reference material; no proprietary profile data is redistributed.
\item \textbf{CR-05.} The half-precision (\texttt{half}) format may not be used
to carry scene-referred linear radiance in any stage prior to the density
transform, because its 11-bit significand is insufficient across the required
$18$-stop working range. See Section~\ref{sec:precision}.
\end{itemize}

\subsection{Acceptance Criteria for the Imaging Model}

The following criteria define ``correct'' for the imaging pipeline and are the
basis of the verification plan in Section~\ref{sec:validation}. They are
stricter than the non-functional requirements because they govern whether the
product's central claim is true.

\begin{enumerate}
\item \textbf{AC-1 (Sensitometric fidelity).} For each shipped negative stock,
the modelled characteristic curve must match the reference $D$--$\logE$ data to
within $\pm 0.03$ density units over the exposure range spanning $\Dmin + 0.1$
to $\Dmax - 0.1$, for all three color records.
\item \textbf{AC-2 (Print chain fidelity).} A synthetic 21-step gray wedge
transported through negative and print must reproduce the reference print
density scale to within $\pm 0.04$ density units and $\pm 0.02$ in the
red-minus-green and blue-minus-green color balance differences.
\item \textbf{AC-3 (Grain spectrum).} The synthesized grain Wiener spectrum
must match the reference stock's measured spectrum to within $\pm 15\%$ of
integrated power over $[0, 50]$ cycles/mm, and RMS granularity must follow the
reference density dependence to within $\pm 10\%$.
\item \textbf{AC-4 (Halation geometry).} For a point source of known
irradiance, the radial falloff of the modelled halo must match the reference
exponential decay length to within $\pm 20\%$ per channel, and the
red:green:blue intensity ratio must match to within $\pm 15\%$.
\item \textbf{AC-5 (Neutrality).} With default parameters and correct white
balance, an 18\% gray card must render neutral to within $\Delta E_{00} \leq
1.5$ against the intended print aim point.
\item \textbf{AC-6 (Monotonicity).} The composite pointwise transfer must be
monotonically non-decreasing in each channel over the full working range for
all reachable parameter combinations. Any parameter combination that violates
this must be unreachable through the UI.
\item \textbf{AC-7 (Reversibility of ordering).} Disabling every optional
stage and setting all stock-independent parameters to identity must produce an
output that differs from a direct working-space-to-output-space conversion by
$\Delta E_{00} \leq 0.5$. This is the pipeline's null test.
\end{enumerate}
